\documentclass[a4paper]{article}
%\usepackage{a4}
%\usepackage{amssymb}
%\usepackage{graphicx}
\begin{document}

\title{The strange formula of Dr. Koide}
\author{Alejandro Rivero\thanks{Zaragoza University at Teruel.  
           {\tt arivero@unizar.es}}}



\maketitle

\begin{abstract}
We do a short historical review of the leptonic mass formula of
Yoshio Koide.
\end{abstract}

\section{Subjective history}
At the end of 1981 Koshio Koide, working in some composite models of
quarks and leptons, had the good or bad fortune of stumbling
over a very simple relationship\cite[f.17]{k1} between the three
charged leptons. This resulted in a prediction of 1.777 GeV for the 
mass of the tau lepton. At that time, such prediction was more than
two sigmas away from the measured value, 1.7842 GeV. So by January of
1983 Koide sends to the Physical Review\cite{k2} a decaffeinated presentation,
holding a purely phenonological point of view, and introducing a corrective
term $\delta$ in order to fit the -then suppossed- experimental value. Still,
this paper paves the way for nascent research on {\it democratic family
mixing}.

This is because Koide mass formula,
$$ (m_e+m_\mu+m_\tau) = \frac 23 (\sqrt{m_e}+\sqrt{m_\mu}+\sqrt{m_\tau})^2$$
can be related to the eigenvectors of the {\it democratic matrix} 
$$\pmatrix{1 1 1 \\ 1 1 1 \\ 1 1 1}$$
As Foot\cite{f} pointed out, the formula is equivalent to ask for
an angle exactly $\pi/4$ between the the eigenvector (1,1,1) and the 
vector formed with the square roots of lepton masses. And the rotation 
around (1,1,1) is of course controlled by taking any basis of the nullspace 
of the democratic mixing matrix. But the multiple uses of this doubly
degenerated matrix are not the theme of this history.

I can only imagine the excitation when some years later the value of the
mass for the tau lepton is revised... and the $delta$ correction parameter
becomes plainly zero: the original prediction was right! Koide revives
his formula and models\cite{k3,k4} but the impact is 
very small. In the Arxiv, R. Foot
\cite{f} suggest the geometrical interpretation, and a couple years
later Esposito and Santorelli\cite{es} revisit it, remarking the
stability of the formula under radiative corrections (below the electroweak
breaking scale, this is). It has not escaped that the down quarks also 
approximate well enough, and some effort is done both by Koide and others
to see how the rest of the particles fit there. Contemporary descents
of this effort come up to neutrinos \cite{newone}.


Meanwhile, the increasing interest on see-saw models gives another
oportunity to the democratic mix, as well to generate a top quark 
enhancement as to justify a radiative generation for the lower generations.
For some years, see-saw marries Koide formula\cite{k5,k6,k7}. The last
incarnation I am aware of, from \cite{k7,k8}, jumps to a justification
based on discrete S3 symmetry. 

Surely, the history is not over yet. After all, as it is remarked in
\cite{newone}, the current status of the fit against a theoretical
unity quotient is $1^{+0.00002635}_{-0.00002021}$. 

\section{More}

%NCG

%dos matrices de masas

%apagando las ctes de acoplo
{\bf Switching couplings off}. It has been remarked and footnoted thousands
of times that the muon and the electon (or equivalently the chyral and
electromagnetic break scales) are separated by a factor of order
alpha. It is less mentioned\cite{y}, but similarly intriguing, that the
tau and the electroweak vacuum (equiv, the SU(3) and electroweak scales) are
separated also by the order of this same Sommerfeld constant. We can mentally
to imagine that if we shut down this constant, the mass of the electron
is pushed towards zero, while the Fermi scale is pushed towards infinity.
In this scenario, Foot's hint of an additional symmetry invites to use
Koide's as a restriction: we find that the net result is that the muon
should increase some MeVs and the tau to decrease a little bit in order
to keep the angle. Moreover, the pion should lost some MeVs of its mass due
to the component quarks, then going to be mass degenerated with the
muon! Moreover, the switching off of the weak coupling would affect to
the link between the bottom quark and the top, causing the former to
loss some weight. In the limit of very small alpha, we have a bunch of
almost massless particles, another one of hugely masive bosons (and the top),
and then some surviving elements making use of the SU(3) gap: muon and 
strange courting the pion, tau charm and bottom dancing around the nucleon
and its glue. The scale between tau and lepton is exactly ${1 + \sqrt 3 
\over 1-sqrt 3}$. This 3 came originally from the number of generations,
not the number of colours. 

One of the puzzling misteries of Nature is why the massive leptons,
colourless, are wandering just there. A hint of universality could
come if we go further, swithing off the SU(3) coupling! Even without
the contribution of strings and glue, we should still could build mesons 
starting at the pion scale, because of the strange quark, and barions at 
the scale of the nucleon, because of the other two massive quarks. Family
masses conspire to save the gap.

%the square roots and Cabibbo angle.

{\bf Cabibbo angle}. 
Since the first publications, Koide's formula has been associated to another
one for the Cabbibo angle, involving square roots of the three generations.
Of course when one the masses is driven to zero, the two extant one
form the very well known square roots of mass quotients that nowadays
are popular folklore in the phenomenology of mixing. If one has followed
the gedakenexercise of the previous parragraph, it is not surprised that
the angle can be obtained both from leptons and quarks. 

\begin{thebibliography}{30}


\bibitem{k1}  {\it A Fermion-Boson Composite Model of quark leptons} 
\bibitem{k2}  {\it New view of quark and lepton mass hierachy } Phys Rev D, Vol 28 p 252 
\bibitem{f} R.  Foot {\it  } hep-ph/9402242 
\bibitem{k3}  {\it Charged Lepton Mass Sum Rule from U(3)-Family Higgs Potential Model} 1989 
\bibitem{es} S. Esposito and P. Santorelli {\it A Geometric Picture for Fermion Masses} hep-ph/9603369 
\bibitem{newone}  {\it Estimate of neutrino masses from Koide's relation } 
\bibitem{k4}  {\it  New Physics from U(3)-Family Nonet Higgs Boson Scenario} 
\bibitem{k5}  {\it  } hep-ph/9505201 
\bibitem{k6}  {\it  } hep-ph/9602303
\bibitem{k7}  {\it Universal Seesaw mass matrix model with an $S_3$ symmetry } Phys Rev D 60 p. 077301 
\bibitem{k8}  {\it  } hep-ph/0005137
\bibitem{y}  Ray J. Yablon, communicated across internet. 
\bibitem{}  {\it  } 
\bibitem{}  {\it  } 
\bibitem{}  {\it  } 
\bibitem{}  {\it  } 
\bibitem{}  {\it  } 
\bibitem{}  {\it  } 
\end{thebibliography}
\end{document}


